\appendix
\section*{Appendix C — Holonomy and Curvature}
\addcontentsline{toc}{section}{Appendix C — Holonomy and Curvature}

\section{Holonomy Framework}

\subsection{Operator-Induced Frame Shifts}
Holonomy measures the frame shift produced by traversing a closed
operator loop. For a loop $\gamma = (O_1, \ldots, O_n)$ with
$\gamma(s_0) = s_0$, the holonomy transformation is
\[
\mathrm{Hol}(\gamma) = O_n \circ \cdots \circ O_1,
\]
acting on the local frame $F$ rather than the state $s$.

\subsection{Holonomy Action on Frames}
The transformed frame is
\[
F' = \mathrm{Hol}(\gamma)(F).
\]
If $F' = F$, curvature is zero. If $F' \neq F$, curvature is present.

\section{Curvature from Non-Commutativity}

\subsection{Commutators}
Curvature arises from operator non–commutativity. For operators
$O_i, O_j$,
\[
[O_i, O_j] = O_j \circ O_i - O_i \circ O_j.
\]

\subsection{Holonomy Tensor}
The holonomy tensor is defined by
\[
H_{ij} = [O_i, O_j],
\]
and encodes curvature contributions from operator pairs.

\subsection{Curvature Accumulation}
For a path $\gamma = (O_1, \ldots, O_n)$,
\[
K(\gamma) = \sum_{i < j} H_{ij},
\]
representing total curvature accumulated along the path.

\section{Definitions}

\begin{definition}[Holonomy]
Holonomy is the frame transformation induced by traversing a closed
operator loop.
\end{definition}

\begin{definition}[Curvature Tensor]
The curvature tensor $H$ is defined by
\[
H_{ij} = [O_i, O_j],
\]
and measures deviation from commutativity.
\end{definition}

\begin{definition}[Path Curvature]
Path curvature is the accumulated holonomy contribution
\[
K(\gamma) = \sum_{i < j} H_{ij}.
\]
\end{definition}

\begin{definition}[Flat Region]
A region is flat if all operators commute, i.e.,
\[
[O_i, O_j] = 0.
\]
\end{definition}

\section{Theorems}

\begin{theorem}[Zero Curvature iff Commutativity]
Curvature is zero if and only if all operators commute.
\end{theorem}

\begin{proof}
Curvature is defined by $H_{ij} = [O_i, O_j]$. If operators commute,
$[O_i, O_j] = 0$ for all pairs, so curvature is zero. Conversely, if any
pair fails to commute, curvature is non–zero.
\end{proof}

\begin{theorem}[Holonomy Non-Reversibility]
Holonomy does not reverse under path inversion:
\[
\mathrm{Hol}(\gamma^{-1}) \neq \mathrm{Hol}(\gamma)^{-1}.
\]
\end{theorem}

\begin{proof}
Holonomy acts on frames, not states. Although
$\gamma^{-1}(\gamma(s)) = s$, the frame satisfies
\[
F' = \mathrm{Hol}(\gamma)(F),
\qquad
F'' = \mathrm{Hol}(\gamma^{-1})(F'),
\]
and in general $F'' \neq F$ because curvature accumulates independently
of state reversibility.
\end{proof}

\begin{theorem}[Curvature Accumulation is Path-Dependent]
For paths $\gamma_1$ and $\gamma_2$,
\[
K(\gamma_2 \circ \gamma_1) = K(\gamma_1) + K(\gamma_2).
\]
\end{theorem}

\begin{proof}
Curvature is defined as a sum over operator pairs. Concatenating paths
adds new operator pairs, producing additive curvature.
\end{proof}

\section{Holonomy and Laplacian Interaction}

\subsection{Curvature-Induced Weight Modification}
Curvature modifies adjacency weights:
\[
w_{ij}' = w_{ij} + H_{ij}.
\]

\subsection{Laplacian Warp}
The Laplacian becomes
\[
L' = D' - W',
\]
producing spectral drift in eigenvalues and eigenvectors.

\section{Discussion}
This appendix formalizes holonomy, curvature tensors, commutators,
curvature accumulation, and their interaction with Laplacian geometry.
These structures form the geometric backbone of narrative physics,
spectral ecology, and reversible quest dynamics.
